\label{p3-20260826:chap:operador-angular-dictamen}
\index{opérateur angulaire}
\index{phase angulaire!sortie spectrale}
\index{clôture!pré-dimensionnelle}

Les deux sections précédentes ont réuni trois éléments qui apparaissaient auparavant
séparés : une identification isométrique entre projecteurs, une mesure de
quotient et un cocycle positif d’action. La présente section étudie la conséquence
spectrale de cette loi et situe le résultat dans l’architecture complète
de HMT\@.

La séparation des statuts est précise. L’identification triple est un théorème
exact relatif à la nouvelle règle variationnelle interne et à un ancrage orienté.
La mesure et le caractère positif sont des théorèmes exacts de \HTr--\HAdd ; le défaut
est en outre relatif à \HDef, la capacité à \HRes et le coût local à \HLoc. La valeur
angulaire est une sortie certifiée par intervalles de l’opérateur \Gtr. Sa
proximité avec la coordonnée angulaire conjuguée canonique est extraordinaire, mais l’intervalle de
la différence exclut l’égalité.

\section{Décomposition de l’opérateur angulaire en \(27\) secteurs de \(72\) états}

Le lecteur de tous les chemins \Gtr possède \(1944\) états microscopiques.
L’action de translation grossière de \((\mathbb Z/3\mathbb Z)^3\) permet la
décomposition
\[
 1944=27\cdot72.
\]
Chaque bloc de \(72\) états est paramétré par
\[
 (r,s,m,d)
 \in(\mathbb Z/3\mathbb Z)^2\times\{0,1\}\times D_4,
\]
où \(m\) est le feuillet et \(d\) la direction précédente. De chaque état
partent quatre continuations cardinales. Le TRIT de la cellule d’arrivée
décide du feuillet suivant.

Avec l’action du \cref{p3-20260826:chap:medida-accion-positiva}, le poids d’une arête
est
\[
 \exp[-g(e)-\lambda_*s(e)].
\]
La transformée en caractères réduit l’opérateur à des matrices complexes
\(L_{k_a,k_b}\) de taille \(72\times72\). Le programme énumère les
arêtes de chaque bloc ; il ne reçoit pas une matrice spectrale précalculée.
L’itération des puissances est conservée à titre exploratoire, mais le résultat qui
suit ne repose plus sur elle : les rayons spectraux sont isolés au moyen d’une
certification rationnelle complète.

Soient \(\rho_{01}\) et \(\rho_{12}\) les rayons spectraux des secteurs
orientés \((0,1)\) et \((1,2)\). Nous définissons la \emph{phase de bande}
\[
 \boxed{y_{\rm band}=\log\rho_{12}-\log\rho_{01}.}
\]
Pour
\[
 \lambda_*
 =1-\frac5{468}-\frac3{11}\Delta_4,
\]
la certification donne
\[
\begin{aligned}
1.460655120862011111
&<\rho_{01}<1.460655120862404024,\\
1.515812008270944328
&<\rho_{12}<1.515812008271195416,
\end{aligned}
\]
et, par propagation par intervalles,
\[
 0.037066226434427529
 <y_{\rm band}<
 0.037066226434862172,
\]
\[
 \boxed{
 2.123738337168943^\circ
 <C_{\rm band}:=\frac{180}{\pi}y_{\rm band}
 <2.123738337193847^\circ.}
\]
La quantité primaire est \(y_{\rm band}\), un logarithme adimensionnel de quotient
spectral. Les degrés apparaissent ensuite lorsque l’on applique la carte
\(180/\pi\). L’opérateur produit d’abord une phase interne et seulement ensuite
une représentation angulaire archimédienne.

\subsection{Calcul exact du rayon spectral de l’opérateur angulaire}

Les blocs sont non normaux. Par conséquent, un petit résidu
\(\|Lv-\mu v\|\) ne suffit pas à démontrer que \(\mu\) est la valeur propre de
plus grand module. La preuve utilisée ici comporte quatre étapes.

On observe d’abord douze lignes structurellement nulles dans chaque bloc. En
développant le polynôme caractéristique suivant ces lignes, on obtient un facteur
\(z^{12}\) et un mineur principal d’ordre soixante. Dans ce dernier apparaissent huit
autres lignes structurellement nulles. Par conséquent,
\[
 \det(zI_{72}-L)=z^{20}\det(zI_{52}-A),
\]
et toute valeur propre non nulle est contenue, avec sa multiplicité, dans un bloc
réduit \(A\) d’ordre cinquante-deux.

Deuxièmement, le certificat fournit des matrices gaussiennes rationnelles
\(S,R\), de dénominateur \(10^{16}\). Le vérificateur calcule exactement
\(E=I-RS\) et obtient
\[
 \|E\|_\infty\le8.5664\times10^{-14}
 \quad\hbox{et}\quad
 \|E\|_\infty\le3.7628\times10^{-14}
\]
pour les secteurs \((0,1)\) et \((1,2)\), respectivement. La série de
Neumann démontre alors que \(S\) est inversible.

Troisièmement, si \(A_0\) est le point médian rationnel du bloc d’intervalles, l’identité
\[
 S^{-1}AS=(I-E)^{-1}RAS
\]
borne la distance à \(C_0=RA_0S\) par
\[
 1.25125\times10^{-13}
 \quad\hbox{et}\quad
 5.7037\times10^{-14}.
\]
Les disques de Gershgorin de \(C_0\), agrandis de ces bornes, contiennent
les disques de la similitude exacte. Dans chaque bloc, un disque est séparé
des cinquante et un autres ; le second théorème de Gershgorin lui attribue
exactement une valeur propre. Son module inférieur dépasse le module supérieur de
tous les autres disques. C’est la preuve que la valeur propre isolée est le
rayon spectral.

Quatrièmement, \(\pi,e,\log\pi,e^{-1},e^{-\lambda_*}\) et \(\sqrt3\) ne sont pas
acceptés comme nombres décimaux de bibliothèque. Le vérificateur les encadre au moyen de
la formule de Machin, des séries factorielle et logarithmique avec reste, de la
série alternée de l’exponentielle et d’un encadrement rationnel quadratique pour
\(\sqrt3\). La vérification emploie uniquement des entiers et des fractions ;
l’argument est clos par la réduction rationnelle, la série de Neumann et
la séparation de Gershgorin exposées dans cette partie. Les développements
décimaux reproductibles à une profondeur de dix mille sont publiés, comme témoins
finis indépendants, dans
\hyperref[app:desarrollos-10000-rev9]{l’annexe des développements}.

\section{Opérateur effectif de bande \texorpdfstring{\(T_{\rm band}\)}{Tband} et isolement de la phase angulaire}

Soit
\[
 x=\frac{50\pi}{9}\alpha_{\rm HMT}
 =0.127362829012869966\ldots
\]
et soit
\[
 R=\begin{pmatrix}1&0\\0&-1\end{pmatrix}.
\]
L’opérateur bicouche effectif associé exclusivement à la phase de bande est
\[
 \boxed{T_{\rm band}=\exp(-xI+y_{\rm band}R).}
\]
Dans la base des secteurs orientés,
\[
 T_{\rm band}
 =\begin{pmatrix}
 e^{-x+y_{\rm band}}&0\\
 0&e^{-x-y_{\rm band}}
 \end{pmatrix}.
\]
On vérifie
\[
 \det T_{\rm band}=e^{-2x},
 \qquad
 \frac{(T_{\rm band})_{11}}{(T_{\rm band})_{22}}
 =e^{2y_{\rm band}}.
\]
Par conséquent,
\[
 -\frac9{100\pi}\log\det T_{\rm band}=\alpha_{\rm HMT},
 \qquad
 \frac{90}{\pi}
 \log\frac{(T_{\rm band})_{11}}{(T_{\rm band})_{22}}
 =C_{\rm band}.
\]

Les deux récupérations ont des statuts distincts. La seconde restitue la
sortie produite par l’écart. La première relit le
\(\alpha_{\rm HMT}\) entré dans \(x\) ; elle démontre la réversibilité de la
carte, et non une seconde dérivation indépendante de \(\alpha\).

La forme exponentielle sépare deux fonctions. Le facteur commun \(e^{-x}\)
décrit la contraction partagée par les deux feuillets. Les facteurs
\(e^{\pm y_{\rm band}}\) distinguent les orientations en conservant le déterminant. Comme
\(I\) et \(R\) commutent, la paire
\[
 \left(\det T_{\rm band},
 \frac{(T_{\rm band})_{11}}{(T_{\rm band})_{22}}\right)
\]
reconstruit univoquement \((x,y_{\rm band})\).

\section{Comparaison ultérieure et échelle de précision}

La section suivante construit une autre phase, la \emph{phase régionale
régularisée} \(y_{\rm reg}\) ---également notée \(y_*\) dans cette section---. Pour
formuler ici une comparaison avec ce résultat ultérieur, nous anticipons la
notation
\[
 C_{\rm reg}:=\frac{180}{\pi}y_{\rm reg}
 =2.123738338968646^\circ.
\]
Cette ligne ne constitue ni une seconde définition ni une entrée de l’opérateur
spectral : elle cite la conclusion du
\cref{p3-20260826:thm:contraangulo-regional}, dont la construction est développée dans la
partie suivante.
Les deux paires d’objets ne sont pas identifiées :
\[
 (y_{\rm band},C_{\rm band})
 \ne
 (y_{\rm reg},C_{\rm reg}).
\]
La différence certifiée par intervalles de la loi linéaire est
\[
 -1.799703\times10^{-9}\ {}^\circ
 <C_{\rm band}-C_{\rm reg}
 <-1.774799\times10^{-9}\ {}^\circ,
\]
avec une différence relative encadrée par
\[
 -8.475\times10^{-10}
 <\frac{C_{\rm band}-C_{\rm reg}}{C_{\rm reg}}
 <-8.356\times10^{-10}.
\]
La loi qui produit \(C_{\rm band}\) ne reçoit ni \(C_{\rm reg}\), ni CKM, ni Barbero,
ni masses ni données SI\@. Cette indépendance à l’égard de la cible fait de la proximité un
candidat de reconnaissance fort et falsifiable. La certification spectrale
élimine l’ancienne lacune numérique ; la différence par intervalles, qui exclut
zéro, fixe désormais la limite exacte : \(C_{\rm band}\) est une sortie spectrale
prédimensionnelle distincte de la réalisation régionale
\(C_{\rm reg}\).

Le terme supplémentaire
\[
 -\frac{\Delta_4^2}{27}
\]
en \(\lambda\) produce
\[
 C_{\rm num}\approx2.12373833894104980{}^\circ,
\]
à \(2.76\times10^{-11}\) degrés de \(C_{\rm reg}\). L’entier \(27\) est interne,
car il compte les caractères de \((\mathbb Z/3\mathbb Z)^3\), mais ce
fait n’impose pas le coefficient du second moment. L’additivité exacte de
\HAdd exclut les termes quadratiques. La variante est conservée comme ablation et non
comme loi principale.

\section{Dépendance fonctionnelle du rayon spectral à l’égard des composantes de l’opérateur}

Toutes les lignes suivantes utilisent le même opérateur ; seule une entrée
de la loi change. Ce sont des diagnostics numériques d’ablation ; l’intervalle certifié
précédent correspond exclusivement à la loi complète
\HTr--\HAdd--\HDef--\HRes--\HLoc (et, en son sein, le défaut positif est fixé dans la
strate \HTr--\HAdd--\HDef).
\[
\begin{array}{l|c|c}
\text{règle}&\lambda&C_{\rm num}\text{ en degrés}\\ \hline
\text{autodual}&1&2.082755794784418\\
\text{microfibre seule}&1-5/468&2.123621809570142\\
\text{trace }5/468+(3/11)\Delta_4
&0.989285913019141&2.123738337181414\\
4/468\text{ au lieu de }5/468
&0.991422665155894&2.115535228138553\\
3/12\text{ au lieu de }3/11
&0.989288440210566&2.123728626433667\\
5/243\text{ au lieu de }5/468
&0.979393542015659&2.161907060464779\\
5/729\text{ au lieu de }5/468
&0.993110963140488&2.109064222037112\\
7/729\text{, mesure des germes}
&0.990367478915522&2.119584299852700.
\end{array}
\]
Les ablations démontrent que la sortie distingue la microfibre complète,
la dimension réduite onze et la mesure des émissions. Elles n’attribuent pas à elles
seules une probabilité au hasard : une telle probabilité exigerait de définir au préalable un
espace de modèles alternatifs et une mesure sur celui-ci.

\section{Composition typée des morphismes archimédien, exceptionnel et dimensionnel}

La construction complète peut s’écrire comme une succession d’applications
aux domaines visibles :
\[
\begin{aligned}
\Omega_{\rm seed}
&\xrightarrow{F}\mathcal U_6
\xrightarrow{J_F}H_U,\\
\mathscr R_\pi^{++}
&\xrightarrow{\beta_\pi}P_GH_G,\\
L_{\rm or}\otimes P_GH_G
&\xrightarrow{W_{GK}}P_3V_{11}
\xrightarrow{U_W}Q_3E_{30},\\
(\Pi_\pi^{(468)},\Delta_4P_3)
&\longmapsto D_{\rm switch}
\xrightarrow{\tau}\delta_*
\xrightarrow{\rm \HRes}\lambda_*,\\
\lambda_*
&\xrightarrow{\rm \HLoc}a_*
\xrightarrow{\mathcal L_{1,*}}y_{\rm band}
\xrightarrow{180/\pi}C_{\rm band},\\
(\alpha_{\rm HMT},y_{\rm band})
&\xrightarrow{\exp(-xI+y_{\rm band}R)}T_{\rm band}.
\end{aligned}
\]
Chaque flèche déclare ce qu’elle conserve :
\begin{itemize}
  \item \(F\) conserve la classe d’émission et laisse dans la fibre sa
  multiplicité de présentation ;
  \item \(J_F\) compense cette multiplicité dans la norme ;
  \item \(\beta_\pi\) conserve la position exceptionnelle ;
  \item \(W_{GK}\) conserve la représentation orientée ;
  \item \(U_W\) conserve l’incidence et la métrique ;
  \item la trace oublie les coordonnées et conserve les rangs ;
  \item l’opérateur de transfert somme les chemins avec le poids de l’action ;
  \item la carte \(180/\pi\) représente une phase interne en degrés.
\end{itemize}

\section{Diagramme des dépendances de l’opérateur angulaire et de ses publications}

\begin{longtable}{>{\raggedright\arraybackslash}p{.31\textwidth}>{\raggedright\arraybackslash}p{.22\textwidth}>{\raggedright\arraybackslash}p{.37\textwidth}}
\toprule
Objet&Nature mathématique&Contenido demostrado\\
\midrule
\endfirsthead
\toprule
Objet&Nature mathématique&Contenido demostrado\\
\midrule
\endhead
\(U_W=I/6:V_{11}\to E_{30}\)
&théorème exact
&isométrie, surjectivité, incidence équivariante et unicité positive\\
\(P_3\leftrightarrow Q_3\)
&théorème exact
&projecteurs conjugués, rang trois et trace \(3/11\)\\
étoile \(R_{B_K}\)
&théorème exact typé
&spectre \(1^{[7]},(-1)^{[4]}\), indice \(3/11\) et obstruction de conjugaison\\
\(P_G\) dans \(t=7\)
&théorème fini exact
&filtration \(11\to3\to1\), projecteur positif et action \(S_3\)\\
\(\mathscr R_\pi^{++}\leftrightarrow P_G\)
&théorème exact
&bijection \(S_3\)-équivariante par position exceptionnelle\\
identification isométrique orientée
&nouveau théorème relatif
&sélecteur unique par \(B_K,u_K\), Reynolds, polaire et réduction \(240\to1\)\\
mesure \(1/468\)
&théorème exact de \HTr
&trace normalisée unique et application polaire induite par les fibres\\
\(D_{\rm switch}\) y \(\delta_*\)
&théorème exact relatif à \HTr--\HAdd--\HDef
&positivité, spectre et trace \(5/468+(3/11)\Delta_4\)\\
\(\lambda_*\)
&théorème interne relatif à \HTr--\HAdd--\HDef--\HRes
&capacité résiduelle obtenue à partir du défaut par la règle \HRes\\
\(\mathcal A_*\)
&théorème interne relatif à \HTr--\HAdd--\HDef--\HRes--\HLoc
&cocycle additif de chemins de coût local \(g+\lambda_*s\)\\
\(T_{\rm band}\) y \(C_{\rm band}\)
&théorème computationnel par intervalles
&réduction \(72\to60\to52\), similitude rationnelle, isolement de
Gershgorin et sortie dont la distance certifiée est comprise entre
\(1.774799\times10^{-9}\) et \(1.799703\times10^{-9}\) degrés de
\(C_{\rm reg}\)\\
\(C_{\rm band}=C_{\rm reg}\)
&affirmation réfutée pour la loi linéaire
&l’intervalle certifié de la différence exclut zéro\\
transition totale
\(\mathcal X_{\mathrm{TPK}}^{\mathrm{enr}}\to \Gtr\)
&réalisateur global requis
&\HTr--\HAdd--\HDef--\HRes--\HLoc définit le lecteur relatif ; la transition
totale exige en outre une mesure image dynamique sur \(\Gtr\)\\
\bottomrule
\end{longtable}

\section{Action prédimensionnelle au centre de la double projection de l’état généalogique}

La double projection part du même état enrichi et produit deux
lectures. Dans la branche archimédienne, les cylindres compatibles se raffinent et
leurs diamètres décroissent ; l’évaluation fixe une valeur limite et oublie une partie
de la généalogie. Dans la branche exceptionnelle, les incidences se prolongent vers
Witt, Golay, Niemeier--Leech, \(V^\natural\), le Monstre et le moonshine.

Le nouveau résultat ajoute un objet commun plus fort qu’une coïncidence de
cardinaux. Le module à onze modes, son secteur positif de rang trois et la
trace de l’action passent de la dodécaphase à Witt par \(U_W\). Le secteur orienté
\(t=7\) s’incorpore par \(W_{GK}\). La microfibre de
\(\pi\) apporte le projecteur régional \(\Pi_\pi^{(468)}\), et la mesure uniforme du
catalogue fixe son poids.

L’action positive se situe ainsi au centre des deux lectures :
\[
 \text{région}
 \quad\oplus\quad
 \text{polarisation orientée}
 \quad\longmapsto\quad
 \text{défaut positif}.
\]
L’évaluation archimédienne et l’évaluation exceptionnelle restent des
applications différentes ; l’entrelaceur identifie un secteur commun, il ne
confond pas leurs codomaines.

\section{Compatibilité de l’opérateur angulaire avec les cylindres du continu discret}

Une histoire HMT détermine des cylindres emboîtés
\[
 C_1\supset C_2\supset\cdots
\]
avec des préfixes compatibles. Lorsque le lecteur archimédien satisfait
\[
 \operatorname{diam}(C_n)\longrightarrow0,
\]
la complétude de \(\mathbb R\) fixe au plus une valeur dans l’intersection.
La nouvelle action ajoute une structure avant cette limite : elle distingue le coût
de la route, la classe d’émission et la multiplicité des présentations que la
valeur finale ne montre plus.

C’est le sens mathématique précis de l’affirmation selon laquelle le continu se lit comme
projection de généalogies discrètes. Le système des cylindres et ses
évaluations se prolongent à toute échelle ; le théorème global de couverture
démontre
\(\operatorname{ev}_{\mathbb R}(\Omega_{\infty,K})=K\) pour chaque compact
\(K\subset\mathbb R\) et, en dirigeant les intervalles, que l’image est tout
\(\mathbb R\). Le théorème fini de cette partie contrôle l’action angulaire
à chaque niveau et s’intègre dans cette couverture globale ; il ne la rouvre ni ne la
conditionne. La descente algébrique de chaque opération conserve sa vérification
propre.

\section{Compatibilité avec le corridor exceptionnel, le Moonshine et les écrans de la théorie M}

L’entrelaceur \(U_W\) aboutit dans le module d’incidence de \(W_{12}\).
À partir du code ternaire, les deux continuations exceptionnelles restent
séparées :
\[
 C_W\longrightarrow W_{12}
 \xrightarrow[\text{relatif à }(D,\ell)]{}W_{24},
\]
et
\[
 C_W\xrightarrow{\ \phi\ } N(A_2^{12})
\xrightarrow{\ v_\alpha\ }\Lambda_{24}
\longrightarrow V_\Lambda
\longrightarrow V^\natural
\longrightarrow\mathbb M
\longrightarrow\text{moonshine}.
\]
L’action positive apporte à cette branche un secteur transporté et une droite
d’orientation. Les flèches de Leech au moonshine sont les constructions
classiques appliquées à l’objet fini obtenu. L’action du Monstre se
réalise sur \(V^\natural\) ; les mots HMT ne sont pas identifiés ici à un
module direct pour cette action.

Les écrans \(12\to11\to10\), \(26\to24\) et l’involution T-duale forment
une interface mathématique exacte avec les dimensions et les dualités utilisées en
théorie M. L’identification physique complète relève de la mécanique
dimensionnelle et est typée dans la seconde partie par l’application supplémentaire vers
\(g_{MN}\), \(C_{MNP}\), \(\psi_M\), l’action et la supersymétrie. Elle n’est pas réservée
à un autre ouvrage et ne se déduit pas de l’action du Monstre.

\section{Factorisation spectrale et isolement des secteurs dominants}

La ligne prédimensionnelle est organisée en quatre strates.
\begin{enumerate}
  \item \textbf{Noyau fini.} APP, TRIT, l’émetteur TPK, le catalogue
  \(104\,976\to468\to243\), les cinq régions de \(\pi\), la
  dodécaphase, \(K\), \(\alpha\), Witt et la branche exceptionnelle possèdent
  des constructions explicites et des certificats.
  \item \textbf{Pont opératoriel.} L’incidence construit \(U_W\) ;
  le secteur orienté \(t=7\) construit \(P_G\) ; \(B_K\) et \(u_K\) sélectionnent
  l’identification isométrique orientée ; les trois projecteurs de rang trois
  sont identifiés relativement à cette règle.
  \item \textbf{Extension d’action.} \HTr fixe la mesure \(1/468\), \HAdd
  fixe le caractère additif et \HDef déclare le couplage
  \(\Delta_4P_3\). \HRes transforme le défaut en capacité et \HLoc porte la
  capacité sur les arêtes. Les cinq clauses produisent la loi locale.
  \item \textbf{Conséquence spectrale.} \Gtr produit
  \(T_{\rm band}\) et une phase angulaire prédimensionnelle certifiée, avec
  une différence relative comprise entre \(8.356\times10^{-10}\) et
  \(8.475\times10^{-10}\) par rapport à \(C_{\rm reg}\).
\end{enumerate}

Cela permet une clôture forte et typée de cette route : l’émetteur fini de la
section coinductive infinie, l’incidence, l’identification des
projecteurs, la mesure et l’action sont clos
relativement aux primitives déclarées, à la règle variationnelle
\(\mathcal E_B\) et aux clauses d’action \HTr--\HAdd--\HDef--\HRes--\HLoc. \HAdd
sélectionne la forme linéaire ; \HDef sélectionne sa seconde coordonnée ; \HRes et \HLoc
précisent comment elle se lit comme coût. Le théorème global de
couverture de tout \(\mathbb R\), l’induction à partir d’une transition totale de
\(\mathcal X_{\mathrm{TPK}}^{\mathrm{enr}}\) et l’équivalence physique avec
la théorie M sont formulés comme des extensions des résultats précédents.

\section{Compatibilité entre les réalisations matricielle, angulaire et réticulaire}

Les représentations construites satisfont simultanément les conditions
de compatibilité suivantes :
\begin{enumerate}
  \item cinq régions de \(\pi\), dont trois positives, une négative et une
  transversale ;
  \item \(P_G,P_3,Q_3\) comme projecteurs positifs et \(R_{B_K}\) comme
  involution d’orientation ;
  \item le secteur homogène orienté \(t=7\), et non le quotient hétérogène de \(t=8\) ;
  \item la droite \(L_{\rm or}\) dans l’identification isométrique triple ;
  \item le caractère relatif de la règle variationnelle qui sélectionne
  \(H_*,c_*,r_*\) ;
  \item \(5/468\) sur les émissions et \(7/729\) sur les germes, sans
  les intervertir ;
  \item la linéarité du caractère comme conséquence de \HAdd et le coefficient
  \(\Delta_4\) comme contenu séparé de \HDef ;
  \item \((y_{\rm band},C_{\rm band})\) comme sortie spectrale certifiée,
  séparée de \((y_{\rm reg},C_{\rm reg})\), la réalisation régionale
  canonique de la section suivante ;
  \item \(\alpha_{\rm HMT}\) comme entrée de \(x\), de sorte que la lecture
  du déterminant est réversible et non indépendante ;
  \item les branches \(W_{12}\to W_{24}\) et
  \(A_2^{12}\to\Lambda_{24}\) comme constructions séparées ;
  \item le lecteur relatif \HTr--\HAdd--\HDef--\HRes--\HLoc distingué d’une
  transition totale, dont la mesure image dynamique constitue une
  construction supplémentaire.
\end{enumerate}

Ces conditions ne sont pas des notes éditoriales. Elles font partie de la structure
logique du résultat.
